\begin{table*}[t]\centering\caption{Main results, $n=10$ planning steps, 20 seeds, mean $\pm$ std. Bold: best, underlined: second best per task and metric. ``Dyna-Q'' is the method under study; the others are reimplemented baselines.}\label{tab:main}
\resizebox{\textwidth}{!}{\begin{tabular}{llcccccccc}
\toprule
Method & Task & cum\_reward $\uparrow$ & post\_reward $\uparrow$ & early\_reward $\uparrow$ & late\_reward $\uparrow$ & first\_goal\_step $\downarrow$ & pre\_reward $\uparrow$ & runtime\_s $\downarrow$ & seeds \\
\midrule
Q-learning & blocking & 5.80 $\pm$ 6.62 & 1.50 $\pm$ 1.67 & 1.55 $\pm$ 1.19 & 0.25 $\pm$ 0.44 & \underline{515.65 $\pm$ 760.07} & 4.30 $\pm$ 7.06 & \textbf{0.00 $\pm$ 0.00} & 20 \\
Dyna-Q+ & blocking & \textbf{92.55 $\pm$ 14.34} & \textbf{51.85 $\pm$ 18.16} & 4.75 $\pm$ 6.54 & \textbf{23.75 $\pm$ 2.63} & \textbf{426.45 $\pm$ 168.97} & \underline{40.70 $\pm$ 13.68} & 0.03 $\pm$ 0.00 & 20 \\
Prioritized sweeping & blocking & \underline{84.25 $\pm$ 38.81} & \underline{29.80 $\pm$ 34.78} & \textbf{18.60 $\pm$ 12.68} & \underline{13.35 $\pm$ 12.81} & 515.65 $\pm$ 760.07 & \textbf{54.45 $\pm$ 27.26} & 0.07 $\pm$ 0.02 & 20 \\
\textbf{Dyna-Q} & blocking & 64.50 $\pm$ 39.90 & 28.25 $\pm$ 37.35 & \underline{7.90 $\pm$ 10.56} & 11.05 $\pm$ 13.71 & 766.90 $\pm$ 788.22 & 36.25 $\pm$ 26.55 & \underline{0.02 $\pm$ 0.00} & 20 \\
\midrule
Q-learning & shortcut & 115.10 $\pm$ 38.94 & 112.00 $\pm$ 38.78 & 0.30 $\pm$ 0.47 & \underline{38.00 $\pm$ 10.55} & \underline{1165.50 $\pm$ 1135.29} & 3.10 $\pm$ 2.00 & \textbf{0.01 $\pm$ 0.00} & 20 \\
Dyna-Q+ & shortcut & \textbf{315.90 $\pm$ 16.82} & \textbf{229.65 $\pm$ 8.29} & 0.05 $\pm$ 0.22 & \textbf{40.85 $\pm$ 2.21} & \textbf{1043.15 $\pm$ 353.15} & 86.25 $\pm$ 14.20 & 0.05 $\pm$ 0.00 & 20 \\
Prioritized sweeping & shortcut & \underline{266.40 $\pm$ 45.53} & \underline{166.15 $\pm$ 16.18} & \textbf{2.65 $\pm$ 5.70} & 28.80 $\pm$ 5.35 & 1165.50 $\pm$ 1135.29 & \textbf{100.25 $\pm$ 51.08} & 0.09 $\pm$ 0.03 & 20 \\
\textbf{Dyna-Q} & shortcut & 256.75 $\pm$ 46.29 & 165.75 $\pm$ 2.00 & \underline{1.80 $\pm$ 4.34} & 27.55 $\pm$ 0.83 & 1194.20 $\pm$ 826.61 & \underline{91.00 $\pm$ 46.49} & \underline{0.04 $\pm$ 0.00} & 20 \\
\midrule
Q-learning & static & 86.25 $\pm$ 55.72 & 75.65 $\pm$ 43.30 & 1.55 $\pm$ 1.19 & 34.40 $\pm$ 13.92 & \underline{457.60 $\pm$ 580.67} & 10.60 $\pm$ 19.11 & \textbf{0.00 $\pm$ 0.00} & 20 \\
Dyna-Q+ & static & \underline{207.80 $\pm$ 12.59} & 124.95 $\pm$ 3.36 & 4.75 $\pm$ 6.54 & 41.25 $\pm$ 3.48 & \textbf{426.45 $\pm$ 168.97} & \underline{82.85 $\pm$ 14.37} & 0.03 $\pm$ 0.00 & 20 \\
Prioritized sweeping & static & \textbf{220.60 $\pm$ 49.98} & \underline{127.75 $\pm$ 14.20} & \textbf{18.60 $\pm$ 12.68} & \underline{44.10 $\pm$ 1.45} & 457.60 $\pm$ 580.67 & \textbf{92.85 $\pm$ 39.97} & 0.04 $\pm$ 0.01 & 20 \\
\textbf{Dyna-Q} & static & 200.15 $\pm$ 48.53 & \textbf{128.40 $\pm$ 17.39} & \underline{7.90 $\pm$ 10.56} & \textbf{44.15 $\pm$ 0.99} & 652.10 $\pm$ 562.45 & 71.75 $\pm$ 38.83 & \underline{0.02 $\pm$ 0.00} & 20 \\
\midrule
Q-learning & stochastic & 49.05 $\pm$ 34.01 & 42.45 $\pm$ 28.32 & 1.10 $\pm$ 1.12 & \textbf{17.30 $\pm$ 10.19} & \underline{501.30 $\pm$ 493.16} & 6.60 $\pm$ 7.03 & \textbf{0.01 $\pm$ 0.00} & 20 \\
Dyna-Q+ & stochastic & 65.60 $\pm$ 10.38 & 36.55 $\pm$ 7.10 & 4.05 $\pm$ 4.39 & 11.25 $\pm$ 5.30 & \textbf{386.75 $\pm$ 207.16} & \textbf{29.05 $\pm$ 8.05} & 0.03 $\pm$ 0.00 & 20 \\
Prioritized sweeping & stochastic & \textbf{74.15 $\pm$ 12.71} & \textbf{48.70 $\pm$ 9.60} & \underline{4.30 $\pm$ 5.03} & \underline{17.05 $\pm$ 4.26} & 501.30 $\pm$ 493.16 & \underline{25.45 $\pm$ 13.24} & 0.16 $\pm$ 0.02 & 20 \\
\textbf{Dyna-Q} & stochastic & \underline{71.05 $\pm$ 28.76} & \underline{46.10 $\pm$ 16.98} & \textbf{4.30 $\pm$ 4.95} & 15.40 $\pm$ 6.13 & 634.75 $\pm$ 678.76 & 24.95 $\pm$ 14.95 & \underline{0.02 $\pm$ 0.00} & 20 \\
\midrule
Q-learning & stoch\_blocking & 5.05 $\pm$ 4.15 & 2.35 $\pm$ 3.60 & 1.10 $\pm$ 1.12 & 0.70 $\pm$ 1.63 & \underline{583.40 $\pm$ 644.51} & 2.70 $\pm$ 2.43 & \textbf{0.01 $\pm$ 0.00} & 20 \\
Dyna-Q+ & stoch\_blocking & \textbf{33.90 $\pm$ 8.72} & \textbf{18.40 $\pm$ 5.28} & 4.05 $\pm$ 4.39 & \textbf{7.70 $\pm$ 3.15} & \textbf{386.75 $\pm$ 207.16} & \textbf{15.50 $\pm$ 6.68} & 0.03 $\pm$ 0.00 & 20 \\
Prioritized sweeping & stoch\_blocking & \underline{22.55 $\pm$ 10.92} & \underline{9.30 $\pm$ 11.87} & \underline{4.30 $\pm$ 5.03} & \underline{4.80 $\pm$ 5.66} & 583.40 $\pm$ 644.51 & 13.25 $\pm$ 10.08 & 0.14 $\pm$ 0.03 & 20 \\
\textbf{Dyna-Q} & stoch\_blocking & 22.00 $\pm$ 15.71 & 8.20 $\pm$ 12.00 & \textbf{4.30 $\pm$ 4.95} & 4.00 $\pm$ 5.66 & 610.00 $\pm$ 647.41 & \underline{13.80 $\pm$ 10.43} & \underline{0.02 $\pm$ 0.00} & 20 \\
\bottomrule
\end{tabular}
}
\end{table*}

\begin{table*}[t]\centering\caption{Cumulative reward of each agent as a function of $n$ (10 seeds, mean). The $n=0$ column is Q-learning on the same seeds.}\label{tab:sweepn}
\resizebox{0.62\textwidth}{!}{\begin{tabular}{llrrrrrr}
\toprule
Task & System & n=0 & n=1 & n=5 & n=20 & n=50 & n=100 \\
\midrule
\multirow{3}{*}{static} & Dyna-Q & \multirow{3}{*}{\rhval{sweep_n/q-learning/static/cum_reward/mean:1}} & \rhval{sweep_n/dyna-q@n=1/static/cum_reward/mean:1} & \rhval{sweep_n/dyna-q@n=5/static/cum_reward/mean:1} & \rhval{sweep_n/dyna-q@n=20/static/cum_reward/mean:1} & \rhval{sweep_n/dyna-q@n=50/static/cum_reward/mean:1} & \rhval{sweep_n/dyna-q@n=100/static/cum_reward/mean:1} \\
 & Dyna-Q+ &  & \rhval{sweep_n/dyna-q+@n=1/static/cum_reward/mean:1} & \rhval{sweep_n/dyna-q+@n=5/static/cum_reward/mean:1} & \rhval{sweep_n/dyna-q+@n=20/static/cum_reward/mean:1} & \rhval{sweep_n/dyna-q+@n=50/static/cum_reward/mean:1} & \rhval{sweep_n/dyna-q+@n=100/static/cum_reward/mean:1} \\
 & PS &  & \rhval{sweep_n/prioritized-sweeping@n=1/static/cum_reward/mean:1} & \rhval{sweep_n/prioritized-sweeping@n=5/static/cum_reward/mean:1} & \rhval{sweep_n/prioritized-sweeping@n=20/static/cum_reward/mean:1} & \rhval{sweep_n/prioritized-sweeping@n=50/static/cum_reward/mean:1} & \rhval{sweep_n/prioritized-sweeping@n=100/static/cum_reward/mean:1} \\
\midrule
\multirow{3}{*}{blocking} & Dyna-Q & \multirow{3}{*}{\rhval{sweep_n/q-learning/blocking/cum_reward/mean:1}} & \rhval{sweep_n/dyna-q@n=1/blocking/cum_reward/mean:1} & \rhval{sweep_n/dyna-q@n=5/blocking/cum_reward/mean:1} & \rhval{sweep_n/dyna-q@n=20/blocking/cum_reward/mean:1} & \rhval{sweep_n/dyna-q@n=50/blocking/cum_reward/mean:1} & \rhval{sweep_n/dyna-q@n=100/blocking/cum_reward/mean:1} \\
 & Dyna-Q+ &  & \rhval{sweep_n/dyna-q+@n=1/blocking/cum_reward/mean:1} & \rhval{sweep_n/dyna-q+@n=5/blocking/cum_reward/mean:1} & \rhval{sweep_n/dyna-q+@n=20/blocking/cum_reward/mean:1} & \rhval{sweep_n/dyna-q+@n=50/blocking/cum_reward/mean:1} & \rhval{sweep_n/dyna-q+@n=100/blocking/cum_reward/mean:1} \\
 & PS &  & \rhval{sweep_n/prioritized-sweeping@n=1/blocking/cum_reward/mean:1} & \rhval{sweep_n/prioritized-sweeping@n=5/blocking/cum_reward/mean:1} & \rhval{sweep_n/prioritized-sweeping@n=20/blocking/cum_reward/mean:1} & \rhval{sweep_n/prioritized-sweeping@n=50/blocking/cum_reward/mean:1} & \rhval{sweep_n/prioritized-sweeping@n=100/blocking/cum_reward/mean:1} \\
\midrule
\multirow{3}{*}{shortcut} & Dyna-Q & \multirow{3}{*}{\rhval{sweep_n/q-learning/shortcut/cum_reward/mean:1}} & \rhval{sweep_n/dyna-q@n=1/shortcut/cum_reward/mean:1} & \rhval{sweep_n/dyna-q@n=5/shortcut/cum_reward/mean:1} & \rhval{sweep_n/dyna-q@n=20/shortcut/cum_reward/mean:1} & \rhval{sweep_n/dyna-q@n=50/shortcut/cum_reward/mean:1} & \rhval{sweep_n/dyna-q@n=100/shortcut/cum_reward/mean:1} \\
 & Dyna-Q+ &  & \rhval{sweep_n/dyna-q+@n=1/shortcut/cum_reward/mean:1} & \rhval{sweep_n/dyna-q+@n=5/shortcut/cum_reward/mean:1} & \rhval{sweep_n/dyna-q+@n=20/shortcut/cum_reward/mean:1} & \rhval{sweep_n/dyna-q+@n=50/shortcut/cum_reward/mean:1} & \rhval{sweep_n/dyna-q+@n=100/shortcut/cum_reward/mean:1} \\
 & PS &  & \rhval{sweep_n/prioritized-sweeping@n=1/shortcut/cum_reward/mean:1} & \rhval{sweep_n/prioritized-sweeping@n=5/shortcut/cum_reward/mean:1} & \rhval{sweep_n/prioritized-sweeping@n=20/shortcut/cum_reward/mean:1} & \rhval{sweep_n/prioritized-sweeping@n=50/shortcut/cum_reward/mean:1} & \rhval{sweep_n/prioritized-sweeping@n=100/shortcut/cum_reward/mean:1} \\
\midrule
\multirow{3}{*}{stochastic} & Dyna-Q & \multirow{3}{*}{\rhval{sweep_n/q-learning/stochastic/cum_reward/mean:1}} & \rhval{sweep_n/dyna-q@n=1/stochastic/cum_reward/mean:1} & \rhval{sweep_n/dyna-q@n=5/stochastic/cum_reward/mean:1} & \rhval{sweep_n/dyna-q@n=20/stochastic/cum_reward/mean:1} & \rhval{sweep_n/dyna-q@n=50/stochastic/cum_reward/mean:1} & \rhval{sweep_n/dyna-q@n=100/stochastic/cum_reward/mean:1} \\
 & Dyna-Q+ &  & \rhval{sweep_n/dyna-q+@n=1/stochastic/cum_reward/mean:1} & \rhval{sweep_n/dyna-q+@n=5/stochastic/cum_reward/mean:1} & \rhval{sweep_n/dyna-q+@n=20/stochastic/cum_reward/mean:1} & \rhval{sweep_n/dyna-q+@n=50/stochastic/cum_reward/mean:1} & \rhval{sweep_n/dyna-q+@n=100/stochastic/cum_reward/mean:1} \\
 & PS &  & \rhval{sweep_n/prioritized-sweeping@n=1/stochastic/cum_reward/mean:1} & \rhval{sweep_n/prioritized-sweeping@n=5/stochastic/cum_reward/mean:1} & \rhval{sweep_n/prioritized-sweeping@n=20/stochastic/cum_reward/mean:1} & \rhval{sweep_n/prioritized-sweeping@n=50/stochastic/cum_reward/mean:1} & \rhval{sweep_n/prioritized-sweeping@n=100/stochastic/cum_reward/mean:1} \\
\bottomrule
\end{tabular}
}
\end{table*}

\begin{figure*}[t]\centering\includegraphics[width=0.95\textwidth]{sweep_n_post.pdf}
\caption{Post-change reward (second half for tasks without a change) versus planning steps $n$; mean $\pm$ s.e.\ over 10 seeds. Dotted: Q-learning.}\label{fig:sweepn}\end{figure*}

\begin{figure}[t]\centering\includegraphics[width=0.85\columnwidth]{curves_blocking.pdf}
\caption{Cumulative reward on the blocking maze (change at step 1000), $n=10$, mean $\pm$ s.e.\ band over 20 seeds; dotted line: change point.}\label{fig:curveblock}\end{figure}

\begin{figure}[t]\centering\includegraphics[width=0.85\columnwidth]{curves_shortcut.pdf}
\caption{Cumulative reward on the shortcut maze (change at step 3000), $n=10$; mean $\pm$ s.e.\ band over 20 seeds; dotted line: change point.}\label{fig:curveshort}\end{figure}

We go hypothesis by hypothesis (Table~\ref{tab:main}, Table~\ref{tab:sweepn}, Figs.~\ref{fig:sweepn}--\ref{fig:curveshort}).

\textbf{H1 (planning helps with a correct model): supported.} On the static maze Dyna-Q at $n=10$ reaches the goal 200.2 times against 86.2 for Q-learning (paired $p=1.6\times10^{-6}$), and the early reward is higher (7.9 versus 1.6, paired $p=0.015$, \texttt{compare\_main\_early\_reward.csv}, but with a standard deviation larger than the mean, so the 500-step gap is only weakly resolved). Cumulative reward keeps rising with $n$ up to 100 for Dyna-Q on this task (Table~\ref{tab:sweepn}). Q-learning is also far below the planners on the blocking and shortcut mazes.

\textbf{H2 (stale model): half supported, half refuted.} The predicted decline of Dyna-Q's post-change reward with $n$ is \emph{not} observed on the blocking maze: it is lowest at $n=1$ and non-monotone afterwards, between roughly 18 and 30 for $n\ge5$ (Fig.~\ref{fig:sweepn}; 10 seeds, large variance). At small $n$ the agent is simply slow, which masks any staleness effect; at the $n$ we tried the model-induced damage is therefore not separable from sampling noise. A failure case is visible in the per-seed spread: Dyna-Q's blocking cumulative reward has a standard deviation of 39.9 around a mean of 64.5, with some seeds reaching the goal rarely after the wall moves. Dyna-Q+ does recover better: post-change reward 51.9 versus 28.2 for Dyna-Q at $n=10$ (paired $p=0.025$, uncorrected), and its advantage grows with $n$ (Fig.~\ref{fig:sweepn}). Fig.~\ref{fig:curveblock} shows the Dyna-Q+ mean curve rising more steeply after the change than Dyna-Q's.

\textbf{H3 (shortcut): supported.} Dyna-Q+ obtains 229.7 post-change goals versus 165.8 for Dyna-Q (paired $p=1.9\times10^{-18}$), and Dyna-Q's value is nearly constant over seeds and $n\ge5$, consistent with it never discovering the new gap (Fig.~\ref{fig:curveshort}); we did not log which path each agent used, so this explanation is an inference from the reward counts. Before the change Dyna-Q+ is not ahead of Dyna-Q (registry pre-change means are close, with Dyna-Q's spread much larger), so its gain is post-change.

\textbf{H4 (stochastic): supported in that planning's benefit shrinks with $n$; not shown to be harmful.} In Table~\ref{tab:sweepn}, Dyna-Q's mean cumulative reward on the stochastic maze falls from 92.4 at $n=1$ to 39.6 at $n=100$ (paired $p=0.0016$ for $n=1$ versus 100). At $n=1$ it beats Q-learning's 55.4 on the same seeds ($p=0.007$); at $n\ge50$ the means (45.8, 39.6) are below Q-learning's but the gaps are within one standard deviation (33.6 for Q-learning; paired $p=0.44$ and $0.14$, 10 seeds, \texttt{tests\_stochastic\_n.csv}), so we cannot say planning is worse than no planning, only that its advantage is gone. Dyna-Q+ and prioritized sweeping show the same decline in means. On the stochastic-blocking maze, both errors combine and all methods earn little (Table~\ref{tab:main}); Dyna-Q+ is best there, but its absolute reward is small.

\textbf{H5 (prioritized sweeping): supported only in the sense of no detectable difference.} At $n=10$ PS and Dyna-Q are not distinguishable on cumulative reward (static 220.6 versus 200.2, $p=0.22$; blocking 84.2 versus 64.5, $p=0.10$; shortcut $p=0.56$) or on post-change reward on blocking (29.8 versus 28.2). With 20 seeds and standard deviations of 40--50 goals, only large effects are detectable, so this is not evidence of equivalence. PS has the best early reward on three tasks (18.6 versus 7.9 for Dyna-Q on blocking and static, paired $p=0.014$), i.e.\ it is faster at the very start, so ``indistinguishable'' holds for cumulative and post-change reward only; its dependence on $n$ is flatter than Dyna-Q's in Table~\ref{tab:sweepn}, which we conjecture is because its priority queue empties (we did not log queue length); it also shares Dyna-Q's failure to find the shortcut.

\textbf{Where the method fails.} Q-learning, Dyna-Q and PS are all poor on the stochastic-blocking maze; Dyna-Q at large $n$ on the stochastic maze has lost its advantage over Q-learning (the means are lower but not significantly so); and Dyna-Q+ at $n=1$ is worse than Dyna-Q on the static maze (122.7 versus 188.8); we conjecture, without a control, that bonus-inflated planning values mislead a model with few samples.
